\chapter*{\textcolor{mycolor}{Abstract}}
\addcontentsline{toc}{chapter}{Abstract}

Today's personal assistants wait to be asked: the user has to type a question, set a reminder or open an app. The
moment help would be useful, such as when leaving home, is precisely the moment nobody thinks of asking for it.

\textbf{LIFEWISE} is a context-aware assistant that observes the situation and acts on its own. A small device
fitted with a camera and sensors (motion, door, light) collects cues about the environment. A lightweight
vision-language model (VLM) describes what the camera sees; a small language model (SLM), preceded by written
rules, matches this description against the context (calendar, habits) and decides whether to alert the user on
their phone. Events are stored in a local SQLite database and can be browsed in a web dashboard. All processing
stays local, on a Raspberry Pi used as an Edge AI test bench.

The prototype targets one pilot scenario: helping the user when leaving home. As the user walks out, LIFEWISE
checks that the phone, keys and wallet are being taken and, if one of them is not visible, invites the user to
check. The alert stays an invitation: a camera can only show that an object is out of view, never that it was
forgotten.

The project is carried out by thirteen engineering students organised in six sub-teams, one per link of the
processing chain. This report covers its context, state of the art, design, implementation and evaluation.

\vspace{1cm}

\noindent\textbf{Keywords}: computer vision, Edge AI, Internet of Things, vision-language model, small language
model, proactive assistance, Raspberry Pi.
